\documentclass[11pt,a4paper]{article}

% ------------------------------------------------------------------
% Paquetes
% ------------------------------------------------------------------
\usepackage[utf8]{inputenc}
\usepackage[T1]{fontenc}
\usepackage[spanish,es-noshorthands]{babel}
\usepackage{lmodern}
\usepackage[margin=2.5cm]{geometry}
\usepackage{xcolor}
\usepackage{listings}
\usepackage{fancyhdr}
\usepackage{enumitem}
\usepackage{hyperref}
\usepackage{tcolorbox}
\usepackage{booktabs}
\usepackage{graphicx}

\hypersetup{
  colorlinks=true,
  linkcolor=azulUAL,
  urlcolor=azulUAL,
  citecolor=azulUAL
}

% ------------------------------------------------------------------
% Colores
% ------------------------------------------------------------------
\definecolor{azulUAL}{RGB}{0,80,140}
\definecolor{grisfondo}{RGB}{245,245,245}
\definecolor{verdecmd}{RGB}{20,110,40}
\definecolor{comentario}{RGB}{110,110,110}
\definecolor{palabraclave}{RGB}{160,20,90}
\definecolor{cadena}{RGB}{170,80,0}

% ------------------------------------------------------------------
% Estilos de listings
% ------------------------------------------------------------------
\lstdefinestyle{base}{
  basicstyle=\ttfamily\small,
  backgroundcolor=\color{grisfondo},
  frame=single,
  framerule=0pt,
  rulecolor=\color{grisfondo},
  breaklines=true,
  breakatwhitespace=true,
  showstringspaces=false,
  columns=fullflexible,
  keepspaces=true,
  xleftmargin=6pt,
  xrightmargin=6pt,
  aboveskip=8pt,
  belowskip=8pt,
}

% Lenguaje YARA (definido a mano: no viene con listings)
\lstdefinelanguage{yara}{
  morekeywords={rule,meta,strings,condition,import,and,or,not,any,all,of,them,
                true,false,ascii,wide,nocase,fullword,filesize,at,in,for},
  sensitive=true,
  morecomment=[l]{//},
  morecomment=[s]{/*}{*/},
  morestring=[b]",
}

\lstdefinestyle{yara}{style=base,language=yara,
  keywordstyle=\color{palabraclave}\bfseries,
  commentstyle=\color{comentario}\itshape,
  stringstyle=\color{cadena}}

\lstdefinestyle{yaml}{style=base,language=,
  keywordstyle=\color{palabraclave}\bfseries,
  commentstyle=\color{comentario}\itshape,
  stringstyle=\color{cadena},
  morekeywords={title,id,status,description,author,date,references,logsource,
                category,product,detection,condition,fields,falsepositives,
                level,tags,selection}}

\lstdefinestyle{consola}{style=base,
  basicstyle=\ttfamily\footnotesize,
  backgroundcolor=\color{black!92},
  commentstyle=\color{white},
  basicstyle=\ttfamily\footnotesize\color{white!95},
}

\lstdefinestyle{c}{style=base,language=C,
  keywordstyle=\color{palabraclave}\bfseries,
  commentstyle=\color{comentario}\itshape,
  stringstyle=\color{cadena}}

% Caja para comandos de terminal
\newtcolorbox{terminal}[1][]{
  colback=black!90,colframe=black!90,
  coltext=white,fontupper=\ttfamily\footnotesize,
  boxrule=0pt,arc=2pt,left=6pt,right=6pt,top=4pt,bottom=4pt,#1}

% Caja de idea clave
\newtcolorbox{ideaclave}{
  colback=azulUAL!6,colframe=azulUAL,
  fonttitle=\bfseries,title=Idea clave,
  boxrule=0.8pt,arc=3pt,left=6pt,right=6pt,top=4pt,bottom=4pt}

% ------------------------------------------------------------------
% Encabezado y pie
% ------------------------------------------------------------------
\pagestyle{fancy}
\fancyhf{}
\lhead{\small YARA y Sigma en la práctica}
\rhead{\small Seguridad en Sistemas}
\cfoot{\thepage}
\renewcommand{\headrulewidth}{0.4pt}

% ------------------------------------------------------------------
% Portada
% ------------------------------------------------------------------
\title{\vspace{-1cm}\textbf{\color{azulUAL} Detección de amenazas con YARA y Sigma}\\[4pt]
\large Un ejemplo práctico para clase: fichero, memoria y logs}
\author{Universidad de Almería \\ \small Material docente}
\date{\today}

\begin{document}

\maketitle
\thispagestyle{fancy}

\begin{abstract}
\noindent
Esta práctica muestra, con ejemplos ejecutables y salidas reales, cómo funcionan
dos de las herramientas de detección más usadas en seguridad defensiva:
\textbf{YARA}, que busca patrones dentro de \emph{ficheros} y de la \emph{memoria}
de procesos en ejecución, y \textbf{Sigma}, un formato genérico de reglas para
\emph{logs} que se traduce al lenguaje de consulta de cada plataforma (Splunk,
Elasticsearch, SQLite, etc.). El hilo conductor es una muestra de laboratorio
inofensiva que cifra sus indicadores en disco y los descifra en memoria, lo que
permite ver con claridad por qué el análisis de RAM es imprescindible.
\end{abstract}

\tableofcontents
\vspace{0.5cm}

% ==================================================================
\section{Objetivos y montaje}
% ==================================================================

Al terminar la práctica el alumnado sabrá:
\begin{itemize}[leftmargin=1.4em]
  \item Escribir reglas YARA y aplicarlas a un \textbf{fichero} en disco.
  \item Aplicar la \textbf{misma} regla a la \textbf{RAM} de un proceso vivo y
        entender por qué el resultado es distinto.
  \item Escribir una regla \textbf{Sigma} y traducirla a varios lenguajes de
        consulta, ejecutándola contra un log real.
\end{itemize}

\begin{ideaclave}
El malware moderno oculta sus indicadores en disco (cifrado, empaquetado,
ofuscación). Solo cuando se ejecuta se descifra en memoria. Por eso una regla que
\emph{falla} contra el fichero puede \emph{acertar} contra la RAM del proceso.
Este es el núcleo pedagógico del ejemplo.
\end{ideaclave}

\subsection*{Herramientas usadas}

\begin{center}
\begin{tabular}{@{}lll@{}}
\toprule
\textbf{Herramienta} & \textbf{Versión} & \textbf{Para qué} \\
\midrule
YARA        & 4.5.8   & Reglas sobre ficheros y memoria \\
yara-python & 4.5.4   & Uso de YARA desde Python (opcional) \\
sigma-cli   & 3.1.0   & Conversión de reglas Sigma \\
SQLite      & 3.45    & Log de ejemplo consultable \\
gcc         & (sistema) & Compilar la muestra de laboratorio \\
\bottomrule
\end{tabular}
\end{center}

\subsection*{Estructura de ficheros del proyecto}

\begin{lstlisting}[style=base]
yara/
  muestras/
    muestra_demo.c        # "malware" de laboratorio (inofensivo)
    compilar.sh           # compila la muestra
    eicar.com             # fichero de prueba estandar EICAR
    informe_limpio.txt    # control negativo (fichero limpio)
  reglas/
    demo_fichero.yar      # reglas para escanear ficheros
    demo_ram.yar          # reglas para escanear memoria
  ejecutar_yara.sh        # script que ejecuta toda la demo de YARA
sigma/
  reglas/
    proceso_sospechoso.yml
    beacon_red.yml
  logs/
    eventos.sql           # log sintetico de creacion de procesos
  ejecutar_sigma.sh       # script que ejecuta toda la demo de Sigma
\end{lstlisting}

% ==================================================================
\section{La muestra de laboratorio}
% ==================================================================

La muestra es un programa en C \textbf{inofensivo}. No se conecta a ningún sitio
ni hace nada dañino: se limita a descifrar en memoria dos indicadores y a esperar
en un bucle. Los indicadores son un dominio de mando y control (C2) ficticio y un
marcador de \emph{payload}. En disco están cifrados con una operación XOR sencilla,
de modo que \textbf{no aparecen en texto plano} en el binario.

\begin{lstlisting}[style=c,caption={Fragmento de \texttt{muestra\_demo.c}},captionpos=b]
#define CLAVE_XOR 0x5A

/* "http://c2.malicioso-demo.local/beacon" cifrado con XOR 0x5A */
static unsigned char c2_cifrado[] = { 0x32,0x2e,0x2e,0x2a, /* ... */ };

static char *descifrar(const unsigned char *src, size_t n) {
    char *out = malloc(n + 1);
    for (size_t i = 0; i < n; i++)
        out[i] = (char)(src[i] ^ CLAVE_XOR);  /* descifrado en RAM */
    out[n] = '\0';
    return out;
}
\end{lstlisting}

\begin{tcolorbox}[colback=orange!8,colframe=orange!70!black,title=Nota de seguridad,
  fonttitle=\bfseries,boxrule=0.6pt,arc=2pt]
La muestra llama a \texttt{prctl(PR\_SET\_PTRACER, PR\_SET\_PTRACER\_ANY, \dots)}
únicamente para que, en la demo, \texttt{yara} pueda leer su memoria sin
privilegios de administrador. En un análisis real se usaría \texttt{sudo} o un
entorno aislado. La muestra no realiza ninguna acción peligrosa.
\end{tcolorbox}

Se compila así, quitando los símbolos para que se parezca a una muestra real:

\begin{terminal}
\$ bash yara/muestras/compilar.sh\\
Compilado: .../yara/muestras/muestra\_demo
\end{terminal}

Comprobamos que las cadenas comprometedoras \textbf{no} están en claro en el binario:

\begin{terminal}
\$ strings yara/muestras/muestra\_demo | grep -Ei "c2|EVIL|malicioso"\\
(sin resultados)
\end{terminal}

% ==================================================================
\section{YARA sobre ficheros}
% ==================================================================

Una regla YARA tiene tres secciones: \texttt{meta} (información descriptiva),
\texttt{strings} (los patrones a buscar) y \texttt{condition} (la expresión
booleana que decide si la regla ``salta''). Los patrones pueden ser texto,
secuencias hexadecimales o expresiones regulares.

\begin{lstlisting}[style=yara,caption={\texttt{yara/reglas/demo\_fichero.yar} (extracto)},captionpos=b]
import "elf"    // inspeccionar cabeceras ELF (Linux)

rule EICAR_Fichero_Prueba : prueba {
    strings:
        $eicar = "EICAR-STANDARD-ANTIVIRUS-TEST-FILE"
    condition:
        $eicar
}

rule Malware_Demo_Disco : demo linux {
    strings:
        $s1 = "[muestra_demo]"
        $s2 = "Cadenas descifradas en memoria"
    condition:
        (elf.type == elf.ET_EXEC or elf.type == elf.ET_DYN)
        and any of ($s*)
}

rule C2_En_Claro : demo c2 {
    strings:
        $dominio = "c2.malicioso-demo.local" ascii wide nocase
        $payload = "EVIL_PAYLOAD_DEMO_v1"     ascii wide
    condition:
        any of them
}
\end{lstlisting}

\subsection*{Ejecución}

La sintaxis es \texttt{yara REGLAS OBJETIVO}. Escaneamos cada muestra:

\begin{terminal}
\$ yara yara/reglas/demo\_fichero.yar yara/muestras/muestra\_demo\\
Malware\_Demo\_Disco yara/muestras/muestra\_demo\\[4pt]
\$ yara yara/reglas/demo\_fichero.yar yara/muestras/eicar.com\\
EICAR\_Fichero\_Prueba yara/muestras/eicar.com\\[4pt]
\$ yara yara/reglas/demo\_fichero.yar yara/muestras/informe\_limpio.txt\\
(sin coincidencias: es el control negativo)
\end{terminal}

\begin{ideaclave}
Fíjate en lo que \textbf{no} ha pasado: la regla \texttt{C2\_En\_Claro} no ha
saltado contra el binario. El dominio del C2 existe en el programa, pero está
cifrado en disco, así que YARA no lo encuentra. Lo veremos aparecer en la sección
siguiente, al escanear la memoria.
\end{ideaclave}

Un detalle didáctico: si escaneamos la carpeta en modo recursivo
(\texttt{yara -r}), la regla \texttt{C2\_En\_Claro} \emph{sí} salta contra el
\textbf{código fuente} \texttt{muestra\_demo.c}, porque ahí el dominio aparece en
un comentario en texto plano. El mismo indicador está oculto en el binario y
visible en el fuente.

\begin{terminal}
\$ yara -r yara/reglas/demo\_fichero.yar yara/muestras/\\
Malware\_Demo\_Disco yara/muestras//muestra\_demo\\
C2\_En\_Claro       yara/muestras//muestra\_demo.c\\
EICAR\_Fichero\_Prueba yara/muestras//eicar.com
\end{terminal}

% ==================================================================
\section{YARA sobre la memoria (RAM)}
% ==================================================================

YARA puede escanear la memoria de un proceso en ejecución pasándole su
\textbf{PID} en lugar de un fichero. Recorre las regiones de memoria del proceso
buscando los mismos patrones. Usamos las reglas de \texttt{demo\_ram.yar}, que
buscan los indicadores \textbf{en claro}.

\begin{lstlisting}[style=yara,caption={\texttt{yara/reglas/demo\_ram.yar}},captionpos=b]
rule Malware_Demo_En_Memoria : demo ram {
    strings:
        $dominio = "c2.malicioso-demo.local" ascii wide nocase
        $payload = "EVIL_PAYLOAD_DEMO_v1"     ascii wide
        $url     = "http://c2.malicioso-demo.local/beacon" ascii
    condition:
        any of them
}

rule Beacon_C2_Generico : ram heuristica {
    strings:
        // URL de beacon http hacia un host 'c2.*'
        $beacon = /https?:\/\/c2\.[a-z0-9.\-]+\/beacon/ ascii nocase
    condition:
        $beacon
}
\end{lstlisting}

\subsection*{Ejecución}

Lanzamos la muestra en segundo plano, anotamos su PID, y comparamos el escaneo
del fichero con el escaneo de la memoria:

\begin{terminal}
\$ ./yara/muestras/muestra\_demo \&\\
{[}muestra\_demo{]} PID 93636\\
{[}muestra\_demo{]} Cadenas descifradas en memoria. Esperando...\\[6pt]
\# 1) La regla de C2 contra el FICHERO en disco: FALLA\\
\$ yara yara/reglas/demo\_ram.yar yara/muestras/muestra\_demo\\
(sin salida)\\[6pt]
\# 2) La MISMA regla contra la RAM del proceso: ACIERTA\\
\$ yara -s yara/reglas/demo\_ram.yar 93636\\
Malware\_Demo\_En\_Memoria 93636\\
0x60278e3202a7:\$dominio: c2.malicioso-demo.local\\
0x60278e3202d0:\$payload: EVIL\_PAYLOAD\_DEMO\_v1\\
0x60278e3202a0:\$url: http://c2.malicioso-demo.local/beacon\\
Beacon\_C2\_Generico 93636\\
0x60278e3202a0:\$beacon: http://c2.malicioso-demo.local/beacon
\end{terminal}

La opción \texttt{-s} muestra, para cada coincidencia, la \textbf{dirección de
memoria} y el \textbf{contenido} encontrado. Vemos los tres indicadores ya
descifrados, exactamente donde el malware los dejó en la RAM.

\begin{ideaclave}
Mismo binario, misma regla, resultado opuesto. En disco los indicadores están
cifrados y la regla falla; en memoria están descifrados y la regla los encuentra.
Esta es la razón por la que el \emph{análisis de memoria} (RAM forensics) es una
técnica central en la respuesta ante incidentes.
\end{ideaclave}

\subsection*{Del PID al volcado de memoria completo}

En un caso real no siempre hay un proceso vivo que escanear. Lo habitual es
capturar un \textbf{volcado de memoria} de toda la máquina y analizarlo en frío,
por ejemplo con \textbf{Volatility}. YARA se integra en ese flujo: se puede
escanear el volcado o los procesos extraídos de él con las mismas reglas. El
principio es idéntico al de esta práctica; solo cambia el origen de los bytes.

% ==================================================================
\section{Sigma: una regla, muchos backends}
% ==================================================================

YARA mira \emph{bytes} (ficheros y memoria). Sigma mira \emph{logs}. Una regla
Sigma se escribe una sola vez en YAML, de forma independiente de la plataforma, y
luego se \textbf{traduce} al lenguaje de consulta concreto: SPL para Splunk,
Lucene o EQL para Elasticsearch, SQL para SQLite, etc. Así una misma lógica de
detección se comparte entre organizaciones que usan herramientas distintas.

\begin{lstlisting}[style=yaml,caption={\texttt{sigma/reglas/proceso\_sospechoso.yml}},captionpos=b]
title: Ejecucion sospechosa con conexion a C2 de demostracion
status: experimental
logsource:
  category: process_creation
  product: linux
detection:
  selection_imagen:
    Image|endswith: '/muestra_demo'
  selection_c2:
    CommandLine|contains: 'c2.malicioso-demo.local'
  condition: selection_imagen or selection_c2
fields: [Image, CommandLine, User]
level: high
tags: [attack.command-and-control, attack.t1071.001]
\end{lstlisting}

\subsection*{Traducción a varios lenguajes}

Con \texttt{sigma convert -t BACKEND} obtenemos la consulta en cada plataforma.
La \textbf{misma} regla produce:

\begin{terminal}
\$ sigma convert -t splunk --without-pipeline proceso\_sospechoso.yml\\
Image="*/muestra\_demo" OR CommandLine="*c2.malicioso-demo.local*"\\
\hspace*{2em}| table Image,CommandLine,User\\[6pt]
\$ sigma convert -t lucene --without-pipeline proceso\_sospechoso.yml\\
Image:*\textbackslash/muestra\_demo OR CommandLine:*c2.malicioso\textbackslash-demo.local*\\[6pt]
\$ sigma convert -t sqlite proceso\_sospechoso.yml\\
SELECT * FROM <TABLE\_NAME> WHERE Image LIKE '\%/muestra\textbackslash\_demo' ESCAPE '\textbackslash'\\
\hspace*{2em}OR CommandLine LIKE '\%c2.malicioso-demo.local\%' ESCAPE '\textbackslash'
\end{terminal}

\subsection*{Ejecución real contra un log}

Para que no sea solo teoría, cargamos un log sintético de creación de procesos en
una base SQLite y ejecutamos la consulta que Sigma ha generado. El log mezcla
actividad normal con dos eventos maliciosos.

\begin{terminal}
\$ sqlite3 eventos.db < eventos.sql\\
\$ CONSULTA=\$(sigma convert -t sqlite proceso\_sospechoso.yml \textbar\ sed 's/<TABLE\_NAME>/eventos/')\\
\$ sqlite3 -header -column eventos.db "\$CONSULTA"
\end{terminal}

El resultado detecta exactamente los dos eventos maliciosos y \textbf{descarta el
ruido} (bash, curl a la web de la universidad, python):

\begin{center}
\small
\begin{tabular}{@{}lll@{}}
\toprule
\textbf{UtcTime} & \textbf{Image} & \textbf{CommandLine} \\
\midrule
2026-09-22 08:02:11 & .../muestra\_demo & ./muestra\_demo \\
2026-09-22 08:02:30 & /usr/bin/curl & curl http://c2.malicioso-demo.local/beacon \\
\bottomrule
\end{tabular}
\end{center}

El primer evento coincide por el \emph{nombre de la imagen}
(\texttt{selection\_imagen}); el segundo, por la \emph{línea de comandos}
(\texttt{selection\_c2}). Los tres eventos benignos no aparecen.

% ==================================================================
\section{YARA y Sigma juntos}
% ==================================================================

Las dos herramientas son complementarias, no rivales. Cubren capas distintas de
la detección:

\begin{center}
\begin{tabular}{@{}p{3.2cm}p{5.2cm}p{5.2cm}@{}}
\toprule
 & \textbf{YARA} & \textbf{Sigma} \\
\midrule
Qué inspecciona & Ficheros y memoria (bytes) & Logs y eventos \\
Dónde actúa & Endpoint, análisis forense, sandbox & SIEM, EDR, plataforma de logs \\
Formato & Reglas \texttt{.yar} propias & YAML genérico traducible \\
Pregunta que responde & ¿Este artefacto contiene el patrón? & ¿Ha ocurrido este evento? \\
\bottomrule
\end{tabular}
\end{center}

\begin{ideaclave}
Un flujo de detección realista los encadena: Sigma alerta en el SIEM de que un
proceso sospechoso se ha ejecutado (evento de log); a raíz de esa alerta, el
analista lanza YARA contra la memoria de ese proceso para confirmar la infección
(bytes). Log $\rightarrow$ alerta $\rightarrow$ confirmación en memoria.
\end{ideaclave}

% ==================================================================
\section{Reproducir la práctica}
% ==================================================================

Todo el ejemplo se ejecuta con dos scripts, desde la raíz del proyecto:

\begin{terminal}
\# Demo completa de YARA (fichero + memoria)\\
\$ bash yara/ejecutar\_yara.sh\\[6pt]
\# Demo completa de Sigma (conversion + consulta real)\\
\$ bash sigma/ejecutar\_sigma.sh
\end{terminal}

\subsection*{Ejercicios propuestos}

\begin{enumerate}[leftmargin=1.6em]
  \item Cambia la clave XOR de la muestra y comprueba que las reglas de RAM siguen
        funcionando sin tocarlas, mientras que buscar el texto cifrado en disco no
        sirve de nada.
  \item Añade a \texttt{demo\_ram.yar} una regla que exija que aparezcan
        \emph{a la vez} el dominio y el payload (\texttt{\$dominio and \$payload}).
  \item Escribe una regla Sigma para la categoría \texttt{network\_connection}
        (ya tienes \texttt{beacon\_red.yml} como modelo) y tradúcela a Splunk.
  \item Convierte la regla Sigma a otro backend disponible
        (\texttt{sigma list targets}) y compara la sintaxis.
\end{enumerate}

\vfill
\begin{center}
\small\color{comentario}
Material docente. Las muestras son de laboratorio y completamente inofensivas.
El dominio \texttt{c2.malicioso-demo.local} es ficticio y no resuelve.
\end{center}

\end{document}
